\section{Translation、Disambiguation 与融合机制}

有了世界模型目标，下一步不是立即挑架构，而是先问多模态究竟提供什么信息。课堂给出一个极有用的二分法：translation（翻译）寻找模态间共享内容，disambiguation（消歧）利用一个模态中额外的信息补足另一个模态。这个区分会直接决定损失函数和融合位置。

\subsection{两种多模态价值：共享事实与互补证据}

翻译式任务假设两个模态描述同一底层事实，因此可以由一边预测另一边；消歧式任务则承认某些事实只存在于其中一个模态，联合观察才足以决定答案。真实系统往往同时包含两者。为了避免把二者混成“融合后效果更好”，下面分别追踪共享信息、模态独有信息以及它们对目标后验分布的不同作用。

\lecturefigure{slide-065-translation.jpg}{Translation：两个模态共享较大信息交集}{00:07:47}

Venn 图中的交集代表可跨模态恢复的内容。例如图像中的“红色车辆”可以对应文本中的同一对象；语音和字幕也可表达近似语义。训练目标通常鼓励对应样本靠近、非对应样本分开，或直接生成另一模态。

\lecturefigure{slide-066-translation-vs-disambiguation.jpg}{Translation 与 disambiguation 的几何差别}{00:08:18}

右图强调现实状态比任一单独观测更大。蓝色和黄色各自覆盖现实的一部分，灰色区域是共享部分；若任务只关注灰色，translation 足够；若答案依赖蓝色或黄色独有区域，就需要 disambiguation。癌症数据尤其如此：H\&E 可见形态，RNA 可见表达程序，两者重合但绝不等价。

\lecturefigure{slide-067-translation-example.jpg}{翻译式多模态学习：跨模态建立统一样本表示}{00:08:50}

图中的生成式人物示例带着一个重要警告：不同模态可能在风格或表面上对齐，却没有捕获底层因果结构。一个模型能从文字生成看似相关的图，并不证明它理解了现实状态；translation 的成功首先说明共享信息可恢复。

\lecturefigure{slide-068-disambiguation.jpg}{Disambiguation：一个模态提供另一个模态没有的事实}{00:08:57}

消歧要求模型允许“非对称证据”。当相机看不到警报声来源时，音频不是图像的翻译，而是增加了新的世界状态约束。数学上可用条件熵说明其价值：若
\[
H(Y\mid X_1,X_2)<H(Y\mid X_1),
\]
则模态 $X_2$ 在给定 $X_1$ 后仍提供关于目标 $Y$ 的额外信息。

\lecturefigure{slide-069-disambiguation-example.jpg}{电影场景示例：看见奔跑并不能决定危险类型}{00:09:43}

只看人物奔跑，可能是火灾、追逐、演习或喜剧；音效、对白和环境线索才让模型更新假设。读这张图时不要把“融合”理解成特征数量增加，而应问：哪个模态让哪些候选解释失去概率质量？这正是后文空间邻域帮助判断细胞状态的逻辑。

\teachervoice{00:07:18--00:09:31，老师反复区分 translation 与 disambiguation：前者强调模态间同一事实，后者强调一边含有另一边没有的关键信息。这个区分是整场课的主轴。}

\begin{warningbox}{一个常见误区：共享 embedding 就等于理解}
对齐损失可以让图文配对在向量空间相近，却可能丢掉对行动关键、但只在单一模态存在的信息。若世界模型只保留交集，它会在最需要互补证据的场景中过度自信。
\end{warningbox}

\subsection{五类融合路线}

确定信息角色后，才进入架构选择。课堂列出五类方案，但明确说明这不是完备分类：它们可以组合，也可能只是在不同层、不同计算预算下实现相似的条件化。为了进行公平比较，不能只看模块名字，而要同时检查融合发生得多早、保留了多少局部细节、缺失模态怎样处理，以及序列长度带来的计算代价。

\lecturefigure{slide-070-fusion-options.jpg}{五类融合：对齐、拼接、cross-attention、额外 token 与条件归一化}{00:10:59}

从早到晚可把这些方法放在一条谱系上：输入级拼接最早，独立编码后对齐最晚；cross-attention 在中间显式交换信息；bonus/CLS token 与 adaptive LayerNorm 则把模态信息压缩为控制信号。融合越早，细粒度交互越充分，但计算成本和缺失模态处理也越复杂。

\lecturefigure{slide-072-contrastive-embedding.jpg}{方案一：对比学习把配对模态拉到共享空间}{00:13:33}

以 CLIP 类目标为例，批内正配对 $(i,i)$ 应比负配对 $(i,j)$ 相似。一个对称损失可写为
\[
\mathcal{L}_{\mathrm{clip}}
=-\frac{1}{2N}\sum_{i=1}^{N}\left[
\log\frac{e^{s_{ii}/\tau}}{\sum_j e^{s_{ij}/\tau}}
+\log\frac{e^{s_{ii}/\tau}}{\sum_j e^{s_{ji}/\tau}}
\right],
\]
其中 $s_{ij}$ 是两个编码器输出的相似度，$\tau$ 是温度。它擅长检索和共享语义，却不自动保留所有局部细节。

\lecturefigure{slide-073-direct-concatenation.jpg}{方案二：在共享空间网格上直接拼接输入}{00:14:30}

RGB 与深度图像逐像素对齐时，拼接非常自然：每个位置直接拥有颜色与距离通道。它的前提是模态具有共同采样单位和可靠 registration（配准）。若 RNA spot、细胞边界和 H\&E patch 的空间坐标不一致，直接拼接会把测量误差误当成生物差异。

\lecturefigure{slide-074-cross-attention-examples.jpg}{方案三：Cross-attention 在两个 token 流之间建立查询}{00:15:35}

图中列出视觉语言和多流 Transformer 的常见用法。Cross-attention 的关键不在“多了一个 attention 模块”，而在角色分工：一个流发出 query，另一个流提供 key/value。这样可以让细胞 token 查询周围组织，或让图像 patch 查询基因表达上下文。

\lecturefigure{slide-076-multihead-self-attention.jpg}{回顾多头自注意力：同一序列内生成 $Q/K/V$}{00:16:01}

Self-attention 对输入 $X$ 计算
\[
Q=XW_Q,\qquad K=XW_K,\qquad V=XW_V,
\]
\[
\operatorname{Attn}(X)=\operatorname{softmax}\!\left(\frac{QK^\top}{\sqrt{d_h}}\right)V.
\]
$d_h$ 是每个 head 的维度。多个 head 允许不同关系并行建模，但所有信息仍来自同一 token 集合。

\lecturefigure{slide-077-cross-attention-mechanism.jpg}{Cross-attention：query 来自一个流，key/value 来自另一个流}{00:16:18}

若 $X$ 是目标流、$Y$ 是条件流，则
\[
Q=XW_Q,\qquad K=YW_K,\qquad V=YW_V.
\]
输出长度跟随 $X$，内容却由 $Y$ 中可检索信息调制。方向不能随意忽略：让 H\&E 查询 RNA 与让 RNA 查询 H\&E，计算成本、输出粒度和缺失模态行为都不同。

\begin{knowledgebox}{Collective fusion 设计问题}
选择融合层时至少要问：token 数量是否相差几个数量级？空间坐标是否一一对应？哪个流需要保留输出分辨率？是否允许某个模态缺失？训练时可见、部署时不可见的模态是否会造成 information leakage？
\end{knowledgebox}

\subsection{Token 条件与 Adaptive LayerNorm}

有时不需要所有 token 两两交互，而只需让一个紧凑条件改变主干网络的计算。额外 token 和 adaptive LayerNorm 都属于这种“控制信号”思路：前者进入注意力序列，后者直接改变归一化后的尺度与偏置。进一步的问题是压缩后的条件是否仍保留稀有局部信号，以及节省的注意力成本是否会以表达能力下降为代价。

\lecturefigure{slide-079-cls-token-fusion.jpg}{方案四：用额外 token 向主干注入模态或类别条件}{00:17:54}

图中一边展示 Vision Transformer 的 CLS token，另一边展示把文本 token、图像 token 或条件 token 拼入序列。额外 token 的优点是复用标准 Transformer；缺点是信息必须先被压缩，且长序列注意力成本可能很高。若一个 token 代表整块肿瘤微环境，它能否保留稀有细胞信号需要实证检验。

\lecturefigure{slide-080-adaptive-layernorm.jpg}{方案五：Adaptive LayerNorm 让条件控制每层特征尺度与偏移}{00:19:24}

LayerNorm 先对隐藏向量标准化，再施加可学习参数。Adaptive LayerNorm 令参数由条件 $c$ 生成：
\[
\operatorname{AdaLN}(h;c)=\gamma(c)\odot
\frac{h-\mu(h)}{\sqrt{\sigma^2(h)+\epsilon}}+\beta(c).
\]
条件网络输出尺度 $\gamma(c)$ 与偏移 $\beta(c)$；它们可来自另一模态、时间步或空间邻域。该方法不增加主序列长度，通常比全量 cross-attention 更省算力，但条件信息会被压缩到调制参数中。

\lecturefigure{slide-082-fusion-summary.jpg}{融合方案回看：没有单一方法垄断所有任务}{00:20:18}

总结页最值得读的是“everything is just token soup”这句半玩笑式判断：很多方法最终都在决定哪些信息变成 token、何时混合、何时压缩。课堂问答进一步指出，adaptive LayerNorm 的条件网络本身也可以包含 cross-attention；因此这些名词不是互斥宗派，而是可组合组件。

\begin{lstlisting}[language=Python,caption={多模态融合器的统一伪代码：显式保留缺失模态与方向}]
def fuse(target_tokens, context_tokens, context_available):
    if not context_available:
        return target_tokens
    context_summary = context_encoder(context_tokens)
    target_tokens = cross_attention(
        queries=target_tokens,
        keys=context_tokens,
        values=context_tokens,
    )
    scale, shift = condition_to_layernorm(context_summary)
    return adaptive_layernorm(target_tokens, scale, shift)
\end{lstlisting}

\teachervoice{00:20:18--00:21:31，课堂提问追问 cross-attention 与 LayerNorm 的关系。老师的回答是：条件分支可以很复杂，甚至先做 cross-attention 再生成缩放与偏移；真正需要比较的是表达能力、参数量与计算效率。}

\begin{warningbox}{架构名称不是实验结论}
“使用 cross-attention”无法说明模型是否学会了正确融合。必须检查缺失模态、错配坐标、单模态捷径、不同患者外推和干预后稳定性。结构图只定义了可能的信息通路，不证明通路被正确使用。
\end{warningbox}

\subsection{本章小结}

Translation 与 disambiguation 是比架构名更稳定的分析工具。对比学习强调共享信息，拼接依赖强配准，cross-attention 允许方向性查询，额外 token 和 adaptive LayerNorm 则把条件压缩成控制信号。实际系统可混合使用它们，但必须围绕任务中的信息缺口与部署约束做选择。

